% figures/fig2-carry-chain.tex — U256 4-limb carry-chain data flow
% Carry arrows placed in the gap BETWEEN the dMIR and x86-64 columns,
% labeled c_0, c_1, c_2 so the chain visually connects the carry-out
% of row N to the carry-in operand of row N+1. The 1:1 lowering
% header conveys that the chain applies at both abstraction levels.
\begin{figure}[t]
\centering
\begin{tikzpicture}[
  every node/.style={font=\footnotesize\ttfamily},
  hdr/.style={font=\small\bfseries\rmfamily},
  link/.style={font=\scriptsize\itshape\rmfamily},
  carry/.style={-{Latex[length=1.8mm]}, semithick, red!70!black},
  carrylab/.style={font=\scriptsize\itshape\rmfamily, red!70!black, inner sep=1pt},
]

% Headings
\node[hdr]  at (1.7, 0.9) {dMIR};
\node[hdr]  at (6.5, 0.9) {x86-64};

% dMIR rows (left column)
\node[anchor=west] (d0) at (0, 0.3)  {r0 = add(a0, b0)};
\node[anchor=west] (d1) at (0, -0.4) {r1 = ADC(a1, b1, c$_0$)};
\node[anchor=west] (d2) at (0, -1.1) {r2 = ADC(a2, b2, c$_1$)};
\node[anchor=west] (d3) at (0, -1.8) {r3 = ADC(a3, b3, c$_2$)};

% x86-64 rows (right column)
\node[anchor=west] (x0) at (5.0, 0.3)  {add rax, rcx};
\node[anchor=west] (x1) at (5.0, -0.4) {adc rbx, rdx};
\node[anchor=west] (x2) at (5.0, -1.1) {adc rsi, rdi};
\node[anchor=west] (x3) at (5.0, -1.8) {adc r8,~~r9};

% "1:1 lowering" link label centered between column headings
\node[link] at (4.0, 0.9) {1:1 lowering};

% Carry chain in MIDDLE column between dMIR and x86 (x = 4.0)
% Arrows from below row N's c-output to above row N+1's c-input;
% c-subscript labels visually adjacent to the arrows.
\draw[carry] (4.0, 0.15)  -- node[carrylab, right=1pt]{c$_0$} (4.0, -0.25);
\draw[carry] (4.0, -0.55) -- node[carrylab, right=1pt]{c$_1$} (4.0, -0.95);
\draw[carry] (4.0, -1.25) -- node[carrylab, right=1pt]{c$_2$} (4.0, -1.65);

\end{tikzpicture}
\caption{U256 4-limb addition lowered to dMIR (left) and x86-64
(right). Each \texttt{ADC} maps 1:1 to one \texttt{adc} instruction;
red arrows in the middle trace the serial carry chain
$c_0, c_1, c_2$, which the lowering preserves and which
instruction-level parallelism cannot hide.}
\label{fig:carry-chain}
\end{figure}
